% =====================================================================
% CAS PCBA Engineer — figures preamble (TikZ schematics)
% Little Henderson · Noetic Compass Press · 2026
%
% PREAMBLE-SAFE: this file is \input into the book preamble by the build.
% It must contain ONLY preamble code (package loads, color defs, \tikzset).
% No \begin{document}, no document-body material.
%
% Visual grammar (shared by every PCBA figure and symbol plate):
%   Domain colors as light fills + dark borders + black text:
%     P (Principled) = gold/amber, C (Cognitive) = blue,
%     E (Evaluative) = red,        O (Operational) = green.
%   Neutral / non-domain blocks = gray.
%   Rounded-rectangle blocks; clean downward signal-flow arrows (-{Latex}).
%   Designed to stay <= ~11cm wide (6x9 measure) and legible in grayscale
%   (rely on borders + labels, not color alone).
% =====================================================================

\usepackage{tikz}
\usetikzlibrary{arrows.meta,positioning,shapes.geometric,calc,fit,backgrounds}

% --- pgfplots for the vector charts (CJT curves, Bode, filters, decay) --
\usepackage{pgfplots}
\pgfplotsset{compat=1.18}
\usepgfplotslibrary{groupplots}            % stacked Bode panels share an x-axis
\pgfplotsset{
  % Shared house style for the analytic charts: light grid, sans labels,
  % measure-friendly default size, grayscale-legible line cycle.
  casaxis/.style={
    width=10.6cm, height=6.2cm,
    font=\sffamily\footnotesize,
    axis line style={draw=casNedge, line width=0.5pt},
    tick label style={font=\sffamily\scriptsize},
    label style={font=\sffamily\footnotesize},
    title style={font=\sffamily\footnotesize\bfseries},
    every axis grid/.style={draw=casNedge!25, line width=0.3pt},
    grid=both, tick align=outside, tickpos=left,
    legend cell align=left,
    legend style={font=\sffamily\scriptsize, draw=casNedge, fill=white,
      inner sep=2pt, row sep=0.5pt}
  }
}

% --- Domain palette (light fills, dark borders) ------------------------
\definecolor{casPfill}{RGB}{255,243,205}   % amber tint  (Principled)
\definecolor{casPedge}{RGB}{153,110,0}     % dark amber
\definecolor{casCfill}{RGB}{213,230,247}   % blue tint   (Cognitive)
\definecolor{casCedge}{RGB}{26,76,140}     % dark blue
\definecolor{casEfill}{RGB}{250,219,216}   % red tint    (Evaluative)
\definecolor{casEedge}{RGB}{150,40,30}     % dark red
\definecolor{casOfill}{RGB}{215,239,221}   % green tint  (Operational)
\definecolor{casOedge}{RGB}{30,110,55}     % dark green
\definecolor{casNfill}{RGB}{233,233,233}   % gray tint   (neutral)
\definecolor{casNedge}{RGB}{90,90,90}      % dark gray

% Accent for "go / pass" and "stop / refuse" paths (still labelled in text,
% so grayscale readers never depend on the color).
\definecolor{casGreen}{RGB}{30,110,55}
\definecolor{casRed}{RGB}{150,40,30}

% --- Reusable block + flow styles --------------------------------------
\tikzset{
  % Base block: rounded rectangle, black text, fixed measure-friendly width.
  casblock/.style={
    rectangle, rounded corners=2pt, draw, line width=0.6pt,
    text=black, align=center, inner sep=3pt,
    minimum height=8mm, font=\sffamily\footnotesize,
    fill=casNfill, draw=casNedge
  },
  % Domain-tinted blocks (inherit casblock, override fill+border).
  pblock/.style={casblock, fill=casPfill, draw=casPedge},
  cblock/.style={casblock, fill=casCfill, draw=casCedge},
  eblock/.style={casblock, fill=casEfill, draw=casEedge},
  oblock/.style={casblock, fill=casOfill, draw=casOedge},
  nblock/.style={casblock, fill=casNfill, draw=casNedge},
  % Wider container block for multi-line stage descriptions.
  widecasblock/.style={casblock, text width=58mm},
  widepblock/.style={pblock, text width=58mm},
  widecblock/.style={cblock, text width=58mm},
  wideeblock/.style={eblock, text width=58mm},
  wideoblock/.style={oblock, text width=58mm},
  widenblock/.style={nblock, text width=58mm},
  % Terminal nodes (signal entry / exit): small pill.
  casterm/.style={casblock, fill=white, draw=black, font=\sffamily\scriptsize\bfseries,
    minimum height=6mm, rounded corners=3pt},
  % Signal-flow arrow (downward by default). The tip carries its own
  % line width=0pt so its apex lands exactly on the path endpoint; without
  % it the tip inherits the wire's 0.6pt stroke and the apex stops short.
  flow/.style={-{Latex[length=2mm,width=1.6mm,line width=0pt]}, line width=0.6pt, draw=black},
  % Feedback / secondary path (dashed).
  fbflow/.style={-{Latex[length=2mm,width=1.6mm,line width=0pt]}, line width=0.5pt, draw=casNedge,
    dashed},
  % Short-hop arrow for links where the free wire is barely 1.5 mm long, so a
  % 2 mm head would overdraw the block it points at.
  flowmini/.style={-{Latex[length=1.2mm,width=1mm,line width=0pt]}, line width=0.6pt, draw=black},
  % Small italic edge label.
  flowlbl/.style={font=\sffamily\scriptsize\itshape, fill=white, inner sep=1pt},
  % Group / section frame drawn behind blocks (backgrounds layer).
  groupframe/.style={rounded corners=3pt, draw=casNedge, line width=0.5pt,
    dash pattern=on 1.5pt off 1.5pt, fill=casNfill, fill opacity=0.25},
  grouplbl/.style={font=\sffamily\scriptsize\bfseries, text=casNedge}
}

% --- Electrical-symbol grammar for Appendix J -------------------------
% These are compact, grayscale-safe functional symbols, not a replacement
% for ANSI/IEC standards. Passive two-terminal pics use local x-terminals at
% (-4mm,0) and (4mm,0). Active, logic, converter, and routing pics declare
% their own leads in the individual definitions; connect to those coordinates.
\tikzset{
  caswire/.style={draw=casNedge, line width=0.6pt},
  cassymbolbox/.style={rounded corners=1.5pt, draw=casNedge, fill=white,
    line width=0.5pt, minimum width=23mm, minimum height=15mm},
  cassymname/.style={font=\sffamily\scriptsize\bfseries, align=center},
  cassymrole/.style={font=\sffamily\tiny, text=casNedge, align=center},
  casnetlbl/.style={font=\sffamily\scriptsize, fill=white, inner sep=1pt},
  casref/.style={font=\sffamily\scriptsize\bfseries},
  casval/.style={font=\sffamily\scriptsize\itshape, text=casNedge},
  casprobe/.style={circle, draw=casNedge, fill=white, line width=0.6pt,
    minimum size=2.6mm, inner sep=0pt},
  % IEC resistor rectangle. Cognitive role: limits activation flow.
  pics/casresistor/.style={code={
    \draw[caswire] (-4mm,0) -- (-2.5mm,0);
    \draw[caswire,fill=white] (-2.5mm,-1mm) rectangle (2.5mm,1mm);
    \draw[caswire] (2.5mm,0) -- (4mm,0);
  }},
  % Non-polar capacitor. Cognitive role: state retention / accumulation.
  pics/cascapacitor/.style={code={
    \draw[caswire] (-4mm,0) -- (-0.55mm,0);
    \draw[caswire,line width=0.9pt] (-0.55mm,-1.7mm) -- (-0.55mm,1.7mm);
    \draw[caswire,line width=0.9pt] (0.55mm,-1.7mm) -- (0.55mm,1.7mm);
    \draw[caswire] (0.55mm,0) -- (4mm,0);
  }},
  % Four-loop inductor. Cognitive role: opposes change in an established flow.
  pics/casinductor/.style={code={
    % Arc endpoints under-report their extrema to PGF; declare the true bounds.
    \useasboundingbox (-4mm,-1mm) rectangle (4mm,1mm);
    \draw[caswire] (-4mm,0) -- (-2.6mm,0);
    \foreach \i in {0,1,2,3}
      \draw[caswire] ({-2.6mm+\i*1.3mm},0) arc[start angle=180,end angle=0,radius=0.65mm];
    \draw[caswire] (2.6mm,0) -- (4mm,0);
  }},
  % Varactor (varicap): a diode whose junction capacitance is controlled
  % by reverse bias. This is not the generic variable-capacitor arrow glyph.
  pics/casvaractor/.style={code={
    \draw[caswire] (-4mm,0) -- (-1.2mm,0);
    \draw[caswire] (1.8mm,0) -- (4mm,0);
    \draw[caswire,fill=white] (-1.2mm,-1.6mm) -- (-1.2mm,1.6mm) -- (0.8mm,0) -- cycle;
    \draw[caswire,line width=0.9pt] (0.8mm,-1.6mm) -- (0.8mm,1.6mm);
    \draw[caswire,line width=0.9pt] (1.8mm,-1.6mm) -- (1.8mm,1.6mm);
  }},
  % Compact diode glyph used by this book's symbol grammar.
  pics/casdiode/.style={code={
    \draw[caswire] (-4mm,0) -- (-1.2mm,0);
    \draw[caswire] (1.2mm,0) -- (4mm,0);
    \draw[caswire,fill=white] (-1.2mm,-1.6mm) -- (-1.2mm,1.6mm) -- (1.2mm,0) -- cycle;
    \draw[caswire,line width=0.9pt] (1.2mm,-1.6mm) -- (1.2mm,1.6mm);
  }},
  % Zener diode: diode with the conventional bent cathode bar.
  pics/caszener/.style={code={
    \draw[caswire] (-4mm,0) -- (-1.2mm,0);
    \draw[caswire] (1.2mm,0) -- (4mm,0);
    \draw[caswire,fill=white] (-1.2mm,-1.6mm) -- (-1.2mm,1.6mm) -- (1.2mm,0) -- cycle;
    \draw[caswire,line width=0.9pt] (1.65mm,-2mm) -- (1.2mm,-1.6mm) -- (1.2mm,1.6mm) -- (0.75mm,2mm);
  }},
  % Coupled coils with a two-line core.
  pics/castransformer/.style={code={
    \foreach \y in {-1.35mm,-0.45mm,0.45mm,1.35mm} {
      \draw[caswire] (-1.15mm,\y) arc[start angle=270,end angle=90,radius=0.45mm];
      \draw[caswire] (1.15mm,\y) arc[start angle=90,end angle=270,radius=0.45mm];
    }
    \draw[caswire] (-0.35mm,-2.2mm) -- (-0.35mm,2.2mm);
    \draw[caswire] (0.35mm,-2.2mm) -- (0.35mm,2.2mm);
  }},
  % Photodiode: diode plus two arrows pointing toward the junction.
  pics/casphotodiode/.style={code={
    \draw[caswire] (-4mm,0) -- (-1.2mm,0);
    \draw[caswire] (1.2mm,0) -- (4mm,0);
    \draw[caswire,fill=white] (-1.2mm,-1.6mm) -- (-1.2mm,1.6mm) -- (1.2mm,0) -- cycle;
    \draw[caswire,line width=0.9pt] (1.2mm,-1.6mm) -- (1.2mm,1.6mm);
    \draw[caswire,-{Latex[length=1mm,width=0.8mm]}] (-2.4mm,3mm) -- (-0.7mm,1.55mm);
    \draw[caswire,-{Latex[length=1mm,width=0.8mm]}] (-0.4mm,3mm) -- (0.5mm,0.75mm);
  }},
  % Three-terminal NPN-style symbol. Appendix J labels the cognitive roles
  % separately because Cause/Bias/Effect are not transistor terminal identities.
  pics/cascjt/.style={code={
    \draw[caswire] (0,0) circle (3.2mm);
    \draw[caswire,line width=0.9pt] (-1.4mm,-2mm) -- (-1.4mm,2mm);
    \draw[caswire] (-4mm,0) -- (-1.4mm,0);
    \draw[caswire] (-1.4mm,1.1mm) -- (0.9mm,2.4mm) -- (0.9mm,4mm);
    \draw[caswire] (-1.4mm,-1.1mm) -- (0.9mm,-2.4mm) -- (0.9mm,-4mm);
    \draw[caswire,-{Latex[length=1mm,width=0.8mm]}] (-0.5mm,-1.6mm) -- (0.7mm,-2.3mm);
  }},
  % Operational amplifier: non-inverting input on top, inverting below.
  pics/casopamp/.style={code={
    \draw[caswire,fill=white] (-3mm,-3mm) -- (-3mm,3mm) -- (3mm,0) -- cycle;
    \draw[caswire] (-5mm,1.6mm) -- (-3mm,1.6mm);
    \draw[caswire] (-5mm,-1.6mm) -- (-3mm,-1.6mm);
    \draw[caswire] (3mm,0) -- (5mm,0);
    \node[font=\sffamily\tiny] at (-2mm,1.6mm) {$+$};
    \node[font=\sffamily\tiny] at (-2mm,-1.6mm) {$-$};
  }},
  % Differential pair: two matched three-terminal devices sharing a tail.
  pics/casdiffpair/.style={code={
    \draw[caswire,fill=white] (-2mm,-1.5mm) circle (1.5mm);
    \draw[caswire,fill=white] (2mm,-1.5mm) circle (1.5mm);
    \draw[caswire] (-4mm,-1.5mm) -- (-3.5mm,-1.5mm);
    \draw[caswire] (4mm,-1.5mm) -- (3.5mm,-1.5mm);
    \draw[caswire] (-2mm,0) -- (-2mm,3mm);
    \draw[caswire] (2mm,0) -- (2mm,3mm);
    \draw[caswire] (-2mm,-3mm) -- (-2mm,-3.6mm) -- (0,-4.4mm);
    \draw[caswire] (2mm,-3mm) -- (2mm,-3.6mm) -- (0,-4.4mm);
  }},
  % Current mirror: two matched branches tied by one bias line. The left
  % device is diode-connected (its base ties to its own collector), and that
  % bias node feeds the right device's base, so the right branch replicates
  % the left branch current. Lead coordinates: collectors up at (+-2, 3.5),
  % emitters on the shared tail at y=-3, bias input at (-3.5, 0).
  pics/casmirror/.style={code={
    \draw[caswire,fill=white] (-2mm,0) circle (1.5mm);
    \draw[caswire,fill=white] (2mm,0) circle (1.5mm);
    \draw[caswire] (-2mm,1.5mm) -- (-2mm,3.5mm);
    \draw[caswire] (2mm,1.5mm) -- (2mm,3.5mm);
    \draw[caswire] (-2mm,-1.5mm) -- (-2mm,-3mm) -- (2mm,-3mm) -- (2mm,-1.5mm);
    % diode connection: left base joined to left collector
    \draw[caswire] (-3.5mm,0) -- (-4.6mm,0) -- (-4.6mm,3.5mm) -- (-2mm,3.5mm);
    % replicated bias line: left base feeds the right base
    \draw[caswire] (-3.5mm,0) -- (-3.5mm,-4.6mm) -- (3.5mm,-4.6mm) -- (3.5mm,0);
  }},
  % Active-load block: composite current-to-difference conversion stage.
  pics/casactiveload/.style={code={
    \draw[caswire,fill=white,rounded corners=1mm] (-3.5mm,-2mm) rectangle (3.5mm,2mm);
    \node[font=\sffamily\tiny\bfseries] at (0,0) {$\Delta I$};
    \draw[caswire] (-5mm,1mm) -- (-3.5mm,1mm);
    \draw[caswire] (-5mm,-1mm) -- (-3.5mm,-1mm);
    \draw[caswire] (3.5mm,0) -- (5mm,0);
  }},
  % Oscillator and crystal timing symbols.
  pics/casoscillator/.style={code={
    \draw[caswire,fill=white] (0,0) circle (3mm);
    \draw[caswire] (-4.5mm,0) -- (-3mm,0); \draw[caswire] (3mm,0) -- (4.5mm,0);
    \draw[caswire] plot[smooth] coordinates {(-2mm,0) (-1.5mm,1.1mm) (-0.5mm,1.1mm) (0.5mm,-1.1mm) (1.5mm,-1.1mm) (2mm,0)};
  }},
  pics/cascrystal/.style={code={
    \draw[caswire] (-4mm,0) -- (-1.6mm,0); \draw[caswire] (1.6mm,0) -- (4mm,0);
    \draw[caswire,line width=0.9pt] (-1.6mm,-1.7mm) -- (-1.6mm,1.7mm);
    \draw[caswire,fill=white] (-0.9mm,-1.3mm) rectangle (0.9mm,1.3mm);
    \draw[caswire,line width=0.9pt] (1.6mm,-1.7mm) -- (1.6mm,1.7mm);
  }},
  pics/caspll/.style={code={
    \draw[caswire,fill=white,rounded corners=0.8mm] (-4mm,-2.6mm) rectangle (4mm,2.6mm);
    \node[font=\sffamily\tiny\bfseries] at (0,0) {PLL};
    \draw[caswire] (-5.5mm,0) -- (-4mm,0); \draw[caswire] (4mm,0) -- (5.5mm,0);
  }},
  % Data converters. Convention: the wide parallel side is the multi-bit
  % port, the narrow side the single-line analog port, and the A/D or D/A
  % text names the conversion direction.
  pics/casadc/.style={code={
    \draw[caswire,fill=white] (-3.5mm,-1.8mm) -- (-3.5mm,1.8mm) -- (3.5mm,2.6mm) -- (3.5mm,-2.6mm) -- cycle;
    \node[font=\sffamily\tiny\bfseries] at (0,0) {A/D};
    \draw[caswire] (-5mm,0) -- (-3.5mm,0); \draw[caswire] (3.5mm,0) -- (5mm,0);
    \draw[caswire] (3.5mm,1.7mm) -- (4.8mm,1.7mm);
    \draw[caswire] (3.5mm,0.9mm) -- (4.8mm,0.9mm);
    \draw[caswire] (3.5mm,-0.9mm) -- (4.8mm,-0.9mm);
    \draw[caswire] (3.5mm,-1.7mm) -- (4.8mm,-1.7mm);
  }},
  pics/casdac/.style={code={
    \draw[caswire,fill=white] (-3.5mm,-2.6mm) -- (-3.5mm,2.6mm) -- (3.5mm,1.8mm) -- (3.5mm,-1.8mm) -- cycle;
    \node[font=\sffamily\tiny\bfseries] at (0,0) {D/A};
    \draw[caswire] (-5mm,0) -- (-3.5mm,0); \draw[caswire] (3.5mm,0) -- (5mm,0);
    \draw[caswire] (-4.8mm,1.7mm) -- (-3.5mm,1.7mm);
    \draw[caswire] (-4.8mm,0.9mm) -- (-3.5mm,0.9mm);
    \draw[caswire] (-4.8mm,-0.9mm) -- (-3.5mm,-0.9mm);
    \draw[caswire] (-4.8mm,-1.7mm) -- (-3.5mm,-1.7mm);
  }},
  % Logic symbols. Small output bubbles produce the inverted variants.
  pics/casand/.style={code={
    \draw[caswire,fill=white] (-3mm,-2.5mm) -- (0,-2.5mm) arc[start angle=-90,end angle=90,radius=2.5mm] -- (-3mm,2.5mm) -- cycle;
    \draw[caswire] (-5mm,1.2mm)--(-3mm,1.2mm); \draw[caswire] (-5mm,-1.2mm)--(-3mm,-1.2mm); \draw[caswire] (2.5mm,0)--(5mm,0);
  }},
  pics/casor/.style={code={
    \draw[caswire,fill=white] (-3mm,-2.5mm) .. controls (-1mm,-1.8mm) and (1.8mm,-1.7mm) .. (3mm,0) .. controls (1.8mm,1.7mm) and (-1mm,1.8mm) .. (-3mm,2.5mm) .. controls (-1.8mm,1mm) and (-1.8mm,-1mm) .. cycle;
    \draw[caswire] (-5mm,1.2mm)--(-2.3mm,1.2mm); \draw[caswire] (-5mm,-1.2mm)--(-2.3mm,-1.2mm); \draw[caswire] (3mm,0)--(5mm,0);
  }},
  pics/casnot/.style={code={
    \draw[caswire,fill=white] (-3mm,-2.5mm)--(-3mm,2.5mm)--(2mm,0)--cycle;
    \draw[caswire] (-5mm,0)--(-3mm,0); \draw[caswire,fill=white] (2.5mm,0) circle (0.5mm); \draw[caswire] (3mm,0)--(5mm,0);
  }},
  pics/casxor/.style={code={
    \draw[caswire] (-3.8mm,-2.5mm) .. controls (-2.6mm,-1mm) and (-2.6mm,1mm) .. (-3.8mm,2.5mm);
    \draw[caswire,fill=white] (-3mm,-2.5mm) .. controls (-1mm,-1.8mm) and (1.8mm,-1.7mm) .. (3mm,0) .. controls (1.8mm,1.7mm) and (-1mm,1.8mm) .. (-3mm,2.5mm) .. controls (-1.8mm,1mm) and (-1.8mm,-1mm) .. cycle;
    \draw[caswire] (-5mm,1.2mm)--(-2.3mm,1.2mm); \draw[caswire] (-5mm,-1.2mm)--(-2.3mm,-1.2mm); \draw[caswire] (3mm,0)--(5mm,0);
  }},
  pics/casnand/.style={code={
    \draw[caswire,fill=white] (-3mm,-2.5mm) -- (0,-2.5mm) arc[start angle=-90,end angle=90,radius=2.5mm] -- (-3mm,2.5mm) -- cycle;
    \draw[caswire] (-5mm,1.2mm)--(-3mm,1.2mm); \draw[caswire] (-5mm,-1.2mm)--(-3mm,-1.2mm); \draw[caswire,fill=white] (3mm,0) circle (0.5mm); \draw[caswire] (3.5mm,0)--(5mm,0);
  }},
  pics/casnor/.style={code={
    \draw[caswire,fill=white] (-3mm,-2.5mm) .. controls (-1mm,-1.8mm) and (1.8mm,-1.7mm) .. (3mm,0) .. controls (1.8mm,1.7mm) and (-1mm,1.8mm) .. (-3mm,2.5mm) .. controls (-1.8mm,1mm) and (-1.8mm,-1mm) .. cycle;
    \draw[caswire] (-5mm,1.2mm)--(-2.3mm,1.2mm); \draw[caswire] (-5mm,-1.2mm)--(-2.3mm,-1.2mm); \draw[caswire,fill=white] (3.5mm,0) circle (0.5mm); \draw[caswire] (4mm,0)--(5mm,0);
  }},
  pics/casxnor/.style={code={
    \draw[caswire] (-3.8mm,-2.5mm) .. controls (-2.6mm,-1mm) and (-2.6mm,1mm) .. (-3.8mm,2.5mm);
    \draw[caswire,fill=white] (-3mm,-2.5mm) .. controls (-1mm,-1.8mm) and (1.8mm,-1.7mm) .. (3mm,0) .. controls (1.8mm,1.7mm) and (-1mm,1.8mm) .. (-3mm,2.5mm) .. controls (-1.8mm,1mm) and (-1.8mm,-1mm) .. cycle;
    \draw[caswire] (-5mm,1.2mm)--(-2.3mm,1.2mm); \draw[caswire] (-5mm,-1.2mm)--(-2.3mm,-1.2mm); \draw[caswire,fill=white] (3.5mm,0) circle (0.5mm); \draw[caswire] (4mm,0)--(5mm,0);
  }},
  % Routing, isolation, hysteresis, regulation, and protection.
  pics/casmux/.style={code={
    \draw[caswire,fill=white] (-3mm,-3mm)--(-3mm,3mm)--(3mm,1.8mm)--(3mm,-1.8mm)--cycle;
    \node[font=\sffamily\tiny\bfseries] at (0,0) {MUX};
    \draw[caswire] (-5mm,1.5mm)--(-3mm,1.5mm); \draw[caswire] (-5mm,-1.5mm)--(-3mm,-1.5mm); \draw[caswire] (3mm,0)--(5mm,0); \draw[caswire] (0,-4mm)--(0,-2.4mm);
  }},
  pics/casdemux/.style={code={
    \draw[caswire,fill=white] (-3mm,-1.8mm)--(-3mm,1.8mm)--(3mm,3mm)--(3mm,-3mm)--cycle;
    \node[font=\sffamily\tiny\bfseries] at (0,0) {DMX};
    \draw[caswire] (-5mm,0)--(-3mm,0); \draw[caswire] (3mm,1.5mm)--(5mm,1.5mm); \draw[caswire] (3mm,-1.5mm)--(5mm,-1.5mm); \draw[caswire] (0,-4mm)--(0,-2.4mm);
  }},
  pics/casrelay/.style={code={
    \draw[caswire] (-4mm,-2mm)--(-2.5mm,-2mm); \draw[caswire,rounded corners=1mm] (-2.5mm,-3mm) rectangle (0.5mm,-1mm); \draw[caswire] (0.5mm,-2mm)--(2mm,-2mm);
    \draw[caswire] (-3mm,2mm)--(-1mm,2mm); \draw[caswire] (-1mm,2mm)--(2mm,3mm); \draw[caswire,fill=white] (2mm,2mm) circle (0.35mm); \draw[caswire] (2mm,2mm)--(4mm,2mm);
    \draw[caswire,dashed] (0,-1mm)--(0,2.33mm);
  }},
  pics/casschmitt/.style={code={
    \draw[caswire,fill=white] (-3mm,-2.5mm)--(-3mm,2.5mm)--(3mm,0)--cycle;
    % Conventional hysteresis mark: two offset threshold steps, kept fully
    % inside the triangle (the lower-right hypotenuse reaches x=0.84 at y=-0.75).
    \draw[caswire,line width=0.45pt] (-1.25mm,-0.75mm)--(-0.35mm,-0.75mm)--(-0.35mm,0.75mm)--(0.45mm,0.75mm);
    \draw[caswire,line width=0.45pt] (-0.85mm,0.75mm)--(0.05mm,0.75mm)--(0.05mm,-0.75mm)--(0.95mm,-0.75mm);
    \draw[caswire] (-5mm,0)--(-3mm,0); \draw[caswire] (3mm,0)--(5mm,0);
  }},
  % Full-wave bridge rectifier: four diode roles around a diamond.
  % Node convention: left and right corners are the AC (~) inputs, the top
  % corner is DC+ and the bottom corner is DC-. Conduction must therefore be
  % upper-left -> DC+, upper-right -> DC+, DC- -> lower-left, DC- -> lower-right.
  % Each rotation below aims that diode's cathode bar at the node it feeds.
  pics/casbridge/.style={code={
    \draw[caswire] (-5mm,0)--(-4mm,0); \draw[caswire] (4mm,0)--(5mm,0);
    \draw[caswire] (0,5mm)--(0,4mm); \draw[caswire] (0,-4mm)--(0,-5mm);
    \draw[caswire] (-4mm,0)--(0,4mm)--(4mm,0)--(0,-4mm)--cycle;
    \foreach \x/\y/\r in {-2/2/45,2/2/135,-2/-2/135,2/-2/45} {
      \begin{scope}[shift={(\x mm,\y mm)},rotate=\r,scale=0.46]
        \draw[caswire,fill=white] (-1.2mm,-1.6mm)--(-1.2mm,1.6mm)--(1.2mm,0)--cycle;
        \draw[caswire,line width=0.9pt] (1.2mm,-1.6mm)--(1.2mm,1.6mm);
      \end{scope}
    }
    \node[font=\sffamily\tiny] at (-5.7mm,0) {$\sim$};
    \node[font=\sffamily\tiny] at (5.7mm,0) {$\sim$};
    \node[font=\sffamily\tiny] at (0,5.8mm) {$+$};
    \node[font=\sffamily\tiny] at (0,-5.8mm) {$-$};
  }},
  pics/casregulator/.style={code={
    \draw[caswire,fill=white,rounded corners=0.8mm] (-3.5mm,-2.5mm) rectangle (3.5mm,2.5mm);
    \node[font=\sffamily\tiny\bfseries] at (0,0) {REG};
    \draw[caswire] (-5mm,0)--(-3.5mm,0); \draw[caswire] (3.5mm,0)--(5mm,0); \draw[caswire] (0,-4mm)--(0,-2.5mm);
  }},
  pics/casfuse/.style={code={
    \draw[caswire] (-4mm,0)--(-2.5mm,0); \draw[caswire,fill=white] (-2.5mm,-1mm) rectangle (2.5mm,1mm); \draw[caswire] (-1.8mm,0)--(1.8mm,0); \draw[caswire] (2.5mm,0)--(4mm,0);
  }},
  pics/casbreaker/.style={code={
    \draw[caswire] (-4mm,0)--(-1.8mm,0); \draw[caswire,fill=white] (-1.8mm,0) circle (0.35mm); \draw[caswire,fill=white] (1.8mm,0) circle (0.35mm); \draw[caswire] (-1.8mm,0)--(1.2mm,1.4mm); \draw[caswire] (1.8mm,0)--(4mm,0);
    % Trip unit: a small square carried on the operating arm (attached, not floating).
    \begin{scope}[shift={(-0.3mm,0.7mm)},rotate=25]
      \draw[caswire,fill=white] (-0.45mm,-0.45mm) rectangle (0.45mm,0.45mm);
    \end{scope}
  }},
  pics/casswitch/.style={code={
    \draw[caswire] (-4mm,0)--(-1.8mm,0); \draw[caswire,fill=white] (-1.8mm,0) circle (0.35mm); \draw[caswire,fill=white] (1.8mm,0) circle (0.35mm); \draw[caswire] (-1.8mm,0)--(1.2mm,1.4mm); \draw[caswire] (1.8mm,0)--(4mm,0);
  }},
  pics/casmeter/.style={code={
    \draw[caswire,fill=white] (0,0) circle (2.6mm); \node[font=\sffamily\tiny\bfseries] at (0,0) {M}; \draw[caswire] (-4mm,0)--(-2.6mm,0); \draw[caswire] (2.6mm,0)--(4mm,0);
  }},
  pics/casrail/.style={code={
    \draw[caswire,line width=1pt] (-4mm,0)--(4mm,0); \draw[caswire] (0,0)--(0,-2.5mm);
  }},
  % Reference/rest marker. It is an engineered convention, not ^0 by identity.
  pics/casreference/.style={code={
    \draw[caswire] (0,0) -- (0,-1.1mm);
    \draw[caswire,line width=0.9pt] (-2.1mm,-1.1mm) -- (2.1mm,-1.1mm);
    \draw[caswire] (-1.35mm,-1.9mm) -- (1.35mm,-1.9mm);
    \draw[caswire] (-0.65mm,-2.7mm) -- (0.65mm,-2.7mm);
  }}
}

% Helper: render a CAS state literal in monospace with carets escaped, in
% math-free text mode. Usage: \casnote{^8^C|A}  ->  prints ^8^C|A safely.
\providecommand{\casnote}[1]{{\ttfamily\scriptsize #1}}
